\section{Arquitectura e integración}
\subsection{Organización de componentes}
La solución combina una aplicación modular con procesos de trabajo especializados y contratos EVM. Los contenedores facilitan su despliegue, pero no justifican describirla como un conjunto de microservicios completamente independientes. API, suscriptor y worker comparten modelos y persistencia, y algunas capas consultan directamente mediante sesiones SQLAlchemy \citep{architecture}.

\begin{figure}[htbp]
\caption{Componentes y relaciones principales del sistema}\label{fig:arquitectura}

\centering
\begin{tikzpicture}[node distance=1cm and .55cm]
\node[box] (iot) {Sensores físicos\\o simulador explícito};
\node[box,right=of iot] (mqtt) {Mosquitto\\y suscriptor};
\node[box,right=of mqtt] (db) {PostgreSQL\\SQLModel / Alembic};
\node[box,below=1.5cm of iot] (web) {React y TypeScript\\Wallet del usuario};
\node[box,right=of web] (api) {FastAPI\\Reglas y permisos};
\node[box,right=of api] (worker) {Worker\\Cálculo e indexación};
\node[box,below=1.5cm of api] (evm) {Contratos EVM\\Registro, SUN y vault};
\node[box,left=of evm] (gov) {Gobernanza\\Timelock};
\node[box,right=of evm] (redis) {Redis\\Revocación y actividad};
\draw[arr] (iot)--(mqtt);
\draw[arr] (mqtt)--(db);
\draw[arr] (web)--(api);
\draw[arr] (api)--(db);
\draw[arr] (db)--(worker);
\draw[arr] (worker)--(evm);
\draw[arr] (gov)--(evm);
\draw[arr] (web.south) to[out=-60,in=160] node[labelnode,above,sloped]{firma del participante} (evm.west);
\draw[arr] (api)--(evm);
\draw[arr] (worker)--(redis);
\end{tikzpicture}
\par\noindent\small\textit{Nota.} Elaboración propia a partir del diseño del repositorio. Las flechas simplifican las relaciones y no representan todas las consultas de lectura \citep{architecture}.

\end{figure}

\subsection{Frontend}
React 19 y TypeScript organizan las pantallas; Vite construye los recursos y nginx sirve la aplicación en Docker. TanStack Query maneja consultas y caché; viem participa en las interacciones EVM; Recharts presenta series de energía. Las cantidades monetarias no deben tratarse como flotantes JavaScript: la API devuelve cadenas enteras de wei y el cliente utiliza conversiones y enteros grandes.

El frontend prepara revisión de red, destinatario, valor y argumentos; la extensión EIP-1193 aporta la firma del usuario. Las claves de los participantes permanecen en su wallet. La sesión HTTP y la cuenta conectada deben coincidir con la relación de wallet del servidor para operar fondos propios.

\subsection{API y autorización}
El punto de entrada es \code{src.main:app}, con prefijo \code{/api/v1}. FastAPI valida solicitudes y publica OpenAPI. Los servicios comprueban propiedad de hogares, acceso comunitario y permisos. El usuario se autentica con contraseñas bcrypt y JWT; la wallet se verifica con un reto EIP-191 de un solo uso.

La API prepara datos de la llamada al contrato sin firmar. La preparación devuelve los datos que el usuario revisa antes de firmar. Los hashes posteriormente recibidos se siguen y contrastan con la cadena. Las rutas y operaciones de la API se resumen en el apartado de resultados \citep{api}.

\subsection{Recepción y trabajo asíncrono}
El suscriptor ejecuta el bucle MQTT y coordina una cola acotada con el procesamiento asíncrono SQL. Confirma mensajes después de la transacción de persistencia. La deduplicación hace que una redelivery por QoS no sume energía dos veces.

El trabajador procesa contadores, finaliza intervalos, descubre lotes, gestiona outbox, transmite operaciones y audita confirmaciones. Redis aporta latidos y revocación de JWT; no es el libro mayor de pagos. El estado de las tareas se conserva para seguimiento y revisión explícita.

\subsection{Registro de transacciones e indexación}
La modificación de negocio y su evento outbox se confirman conjuntamente. Antes de transmitir una operación del oráculo se conservan sus bytes firmados, hash, nonce y origen. Un bloqueo asesor SQL coordina nonces por chain ID, génesis y remitente. Reintentar una respuesta ambigua no debe crear otro pago con un nonce nuevo.

Una operación se refleja como confirmada después de comprobar recibo exitoso, bloque canónico, profundidad y eventos esperados. Las reorganizaciones y los eventos huérfanos exigen reconstrucción de proyecciones o revisión; la aplicación no inventa una reversión económica de lo ya pagado \citep{architecture,blockchain}.

\subsection{Responsabilidades de los componentes}
Los dispositivos aportan las lecturas. El servidor comprueba su origen y calcula la energía. La gobernanza aprueba el lote y los contratos verifican las pruebas y las condiciones económicas. PostgreSQL conserva la evidencia y la aplicación consulta los resultados confirmados de la cadena.
